% Nome del file: ManualeSviluppatore.tex
% Percorso: \gl{template}
% Autore: Vault-Tech
% Data creazione: 10.05.2016
% E-mail: vaulttech.swe@gmail.comcom
%
% Diario delle modifiche: interno al file.

\documentclass[a4paper, titlepage]{article}

\usepackage[margin=3cm]{geometry}
\usepackage{../../Stile}
\usepackage{../../Comandi}

\setcounter{secnumdepth}{5}
\setcounter{tocdepth}{5}

\def\NOME{Manuale Sviluppatore}
\def\VERSIONE{2.0}
\def\DATA{06.06.2016}
\def\REDATTORE{Rudy Berton \\ & Filippo Tesser}
\def\VERIFICATORE{Michela De Bortoli}
\def\RESPONSABILE{Simone Boccato}
\def\USO{Esterno}
\def\DISTRIBUZIONE{\COMMITTENTE \\ & \CARDIN \\ & \PROPONENTE}


\begin{document}
	
	\pagestyle{fancy}	
	\pagenumbering{Roman}
	\rfoot{Pagina \thepage{} di \pageref{lastromanpage}}
	
	\maketitle
	
	\begin{diario}
	\recap{Approvazione del documento}{Michela De Bortoli}{Responsabile}{06.04.2016}{3.0}
	\recap{Correzione errori individuati}{Michela De Bortoli}{Analista}{06.04.2016}{2.10}
	\recap{Verifica dell'intero documento}{Rudy Berton}{Verificatore}{05.04.2016}{2.9}
	\recap{Stesura appendice D}{Giacomo Beltrame}{Analista}{04.04.2016}{2.8}
	\recap{Verifica appendici A e B}{Giacomo Beltrame}{Verificatore}{03.04.2016}{2.7}
	\recap{Stesura test di integrazione}{Rudy Berton}{Amministratore}{02.04.2016}{2.6}
	\recap{Stesura test di sistema}{Vassilikì Menarin}{Progettista}{02.04.2016}{2.5}
	\recap{Modifica della sezione A.3.3 dell'appendice}{Filippo Tesser}{Analista}{01.04.2016}{2.4}
	\recap{Incremento test di accettazione}{Michela De Bortoli}{Progettista}{01.04.2016}{2.3}
	\recap{Inizio stesura specifica dei test (appendice B)}{Michela De Bortoli}{Progettista}{31.03.2016}{2.2}
	\recap{Incremento dell'appendice A}{Filippo Tesser}{Analista}{31.03.2016}{2.1}
	\recap{Approvazione documento}{Miki Violetto}{Responsabile}{23.02.2016}{2.0}
	\recap{Verifica delle sezioni modificate}{Rudy Berton}{Verificatore}{22.02.2016}{1.2}
	\recap{Revisione correttiva dei contenuti rispetto alle segnalazioni del committente}{Giacomo Beltrame}{Analista}{20.02.2016}{1.1}
	\recap{Approvazione documento}{Vassilikì Menarin}{Responsabile}{20.01.2016}{1.0}
	\recap{Verifica del documento}{Simone Boccato}{Verificatore}{19.01.2016}{0.9}
	\recap{Stesura appendice D}{Rudy Berton}{Analista}{18.01.2016}{0.8}
	\recap{Correzione errori segnalati}{Rudy Berton}{Analista}{16.01.2016}{0.7}
	\recap{Verifica del documento}{Filippo Tesser}{Verificatore}{15.01.2016}{0.6}
	\recap{Stesura appendici A, B e C}{Rudy Berton}{Analista}{11.01.2016}{0.5}
	\recap{Fine stesura Gestione della qualità e stesura sezione Gestione amministrativa della revisione}{Rudy Berton}{Analista}{08.01.2016}{0.4}
	\recap{Inizio stesura Gestione della qualità}{Rudy Berton}{Analista}{05.01.2016}{0.3}
	\recap{Stesura sezione Obiettivi di qualità}{Rudy Berton}{Analista}{03.01.2016}{0.2}
	\recap{Stesura sezione Introduzione}{Rudy Berton}{Analista}{02.01.2016}{0.1}
\end{diario}
	
	\newpage
	\tableofcontents\label{lastromanpage}
%	\newpage
%	\listoffigures
%	\newpage
%	\listoftables
	
	\newpage
	\clearpage	
	\pagenumbering{arabic}
	\rfoot{Pagina \thepage{} di \pageref*{LastPage}}
	%Deve esserci per permettere i riferimenti incrociati di colore blu
	\hypersetup{linkcolor=blue}
	
	\section{Introduzione}
	\subsection{Scopo del documento}
	Questo documento ha lo scopo di indicare e spiegare quali sono i comandi da eseguire per installare correttamente l'applicativo Quizzipedia, come creare nuove tipologie di domande ed integrarle nel sistema e come amministrare gli enti all'interno del proprio \gl{server}.
	
	\subsection{Scopo del prodotto}
	\SCOPO
	
	\subsection{Glossario}
	Al fine di evitare ogni ambiguità nel linguaggio e massimizzare la comprensione dei documenti, i termini tecnici, gli acronimi e le abbreviazioni che necessitano di definizione sono riportati nell'\hyperref[gl]{appendice A}.
	Inoltre ogni occorrenza di un vocabolo presente nel “Glossario” sarà posta in corsivo e seguita da
	una ‘g’ maiuscola a pedice (p.es. \gl{Parola}).
		
	\subsection{Riferimenti}	
	\subsubsection{Riferimenti informativi}
	\begin{itemize}
		\item \textbf{Manuale utente \gl{Git}:} \url{https://git-scm.com/docs/user-manual.html};
		\item \textbf{Manuale installazione \gl{Node.js}:} \url{https://docs.npmjs.com/getting-started/installing-node};
		\item \textbf{Manuale installazione \gl{MongoDB}:} \url{https://docs.mongodb.com/master/installation/};
		\item \textbf{Front-end package manager Bower}: \url{http://bower.io/};
		\item \textbf{Grunt JavaScript task runner}: \url{http://gruntjs.com/}.
	\end{itemize}
	\newpage
	
	\section{Prerequisiti software}
	Per poter avviare Quizzipedia è richiesta l'installazione dei seguenti \gl{software}.
	
	\subsection{Sistema operativo}
	Per poter avviare il \gl{server} è richiesta una versione Linux \gl{LTS} a 64bit. \\È necessario disporre dei privilegi di superuser per installare i vari software necessari.
	
	
	\subsection{Node.js e npm}
	Scaricare da browser e installare \gl{Node.js} e il package manager \gl{npm} da\\ \url{https://nodejs.org/en/download/current/}.\\
	\\Oppure da riga di comando con il comando\\ \texttt{\$ wget https://nodejs.org/download/release/v6.0.0/node-v6.0.0-linux-x64.tar.xz} \\e \\ \texttt{\$ tar -C /usr/local --strip-components 1 -xJf node-v6.0.0-linux-x64.tar.xz} \\(dopo essersi posizionati nella cartella dove è stato scaricato ).
	
	Verificare che sia visibile la versione installata con i seguenti comandi:
	
	\texttt{\$ node -v}
	
	\texttt{\$ npm -v}
	
	In caso di problemi nell'installazione si consultino le guide nei riferimenti.

	\subsection{MongoDB}
	Scaricare e installare \gl{MongoDB} da\\ \url{https://www.mongodb.com/download-center}.\\
	\\Oppure da riga di comando con il comando\\
	\texttt{\$ apt-get install mongodb}
	
	Per avviare \gl{MongoDB} sarà necessario dare il seguente comando:
	
	\texttt{\$ mongo}
	
	In caso di problemi nell'installazione si consultino le guide nei riferimenti.

	\subsection{Forever}
	'Forever' è un modulo \gl{npm} che permette l'esecuzione continua di uno script \gl{JavaScript}.
	Scaricare e installare globalmente questo modulo utilizzando \gl{npm} con il seguente comando:

	\texttt{\$ npm install -g forever}
	
	\subsection{Angular.js}
	Scaricare e installare \gl{Angular.js} utilizzando \gl{npm} con il seguente comando:

	\texttt{\$ npm install angular}
	
	\section{Requisiti hardware}
	Il prodotto lato \gl{server} richiede come requisiti minimi:
	\begin{itemize}
		\item processore dual core;
		\item 2GB di memoria RAM;
		\item 1GB di spazio libero su disco.
	\end{itemize}
	
	\section{Installazione e Avvio}
	
	\subsection{Ottenere Quizzipedia}
	\subsubsection{Download da browser}
	Reperire il progetto Quizzipedia è possibile recandosi con qualunque \gl{browser} all'indirizzo \url{http://github.com/vault-tech/Quizzipedia}.
	A questo punto è possibile scaricare il codice sorgente con l'apposito tasto \texttt{Download ZIP}.

	\subsubsection{Download da riga di comando}
	Dopo aver aperto un terminale eseguire il seguente comando:

	\texttt{\$ git clone http://github.com/vault-tech/Quizzipedia dest}

	dove "dest" va sostituita con il percorso alla cartella nella quale si vuole installare Quizzipedia.

	\subsection{Installazione dipendenze}
	
	Una volta ottenuto il codice sorgente, spostarsi all'interno della root del \gl{repository}, si controlli che sia presente il file package.json.
	
	Sempre nella directory corrente eseguire poi il seguente comando da terminale:
	
	\texttt{\$ npm install}
	
	in questo modo verranno installati automaticamente i moduli di \gl{Node.js} necessari per l'esecuzione del prodotto.
	
	\subsection{Configurazione}
	Nella cartella \texttt{config} si trovano i seguenti file di configurazione:
	\begin{itemize}
		\item "email.js": contiene la email che verrà utilizzata per contattare il gestore del prodotto;
		\item "database.js": contiene l'indirizzo del database a cui si connetterà l'applicativo, di default viene usato un database in locale;
		\item "upload.js": contiene gli indirizzi delle cartelle che verranno usate per collezionare gli allegati inviati al \gl{server};
		\item "server.js": contiene la porta sulla quale si vuole rendere reperibile Quizzipedia via \gl{browser} e la durata della sessione,
		se viene specificata una porta inferiore a 1024 allora sarà necessario anteporre \texttt{sudo} al comando di avvio per autorizzare \gl{Node.js} ad utilizzare tale porta.
	\end{itemize}
	La modifica di questi valori è desiderabile ma non obbligatoria al funzionamento del prodotto.
	
	\subsection{Avvio}
	
	Per avviare Quizzipedia e mantenerlo in esecuzione, utilizzare i seguenti comandi da terminale dopo essersi spostati nella root di Quizzipedia:
	\begin{itemize}
		\item \texttt{\$ mongo} per avviare \gl{mongoDB} nel caso venga utilizzato un database in locale,
		\item \texttt{\$ npm start} per avviare il \gl{server} senza utilizzare il modulo 'forever',
		\item \texttt{\$ npm run foreverStart} per avviare il \gl{server} utilizzando il modulo 'forever'.
	\end{itemize}
	Per terminare l'esecuzione di un \gl{server} avviato con 'forever' si utilizzi il comando \texttt{\$ npm run foreverStop}.
	
	\newpage
	
	\section{Creazione nuova tipologia di domande}
	
	\subsection{Sintassi QML}
	\label{qml}
	Il QML (Quiz Markup Language) è un linguaggio con una sintassi definita da utilizzare per il salvataggio delle domande utilizzate nei quiz.
	Vengono qui illustrate le regole da adottare per utilizzare il QML per definire nuove tipologie di domande:
	\begin{itemize}
		\item deve iniziare e finire con i caratteri "\textbar";
		\item la tipologia di domanda è una stringa lunga a piacere ed univoca, deve essere preceduta dalla sequenza di caratteri "q?";
		\item la parte relativa al testo della domanda è una stringa lunga a piacere, deve essere preceduta dalla sequenza di caratteri "\#t\#";
		\item la parte relativa a/alle risposta/e è una stringa lunga a piacere, deve essere preceduta dalla sequenza di caratteri "\#a\#";
		\item nel caso le due suddette stringhe siano multiple, è possibile utilizzare il carattere "§" per separare la varie parti;
		\item per aggiungere una validità o un id a una parte di testo si utilizza una stringa racchiusi  nei caratteri "[" e "]";
		\item eventuali allegati vanno racchiusi tra i caratteri "\{" e "\}", si veda la sottosezione seguente per i dettagli;
		\item alla fine delle stringhe relative al testo e alle risposte, deve essere utilizzata la sequenza di caratteri "\#££\#";
		\item è possibile inserire nuove parti aggiuntive nella stringa QML qualora vengano rispettate queste condizioni e il QMLAgent venga esteso in modo opportuno.
	\end{itemize}
	Segue un esempio di QML per i tipi di domande Vero/Falso e Risposta Multipla.
	\begin{verbatim}|q?trfs#t#Parigi è in Germania#a#false#££#|\end{verbatim}
	\begin{verbatim}|q?mult#t#Quali sono città europee?#a#Londra[true]§Mosca[false]#££#|\end{verbatim}
	
	\subsubsection{Allegati}
	Gli allegati possono essere aggiunti sia per il testo della domanda, che per ogni risposta.
	Si possono utilizzare immagini, video o audio, con un limite di 20MB.
	La sintassi specifica per allegare un file a una domanda o risposta è lo stesso ed è il seguente:
	\begin{itemize}
		\item deve iniziare e finire con i caratteri "\{" e "\}";
		\item deve essere preceduto dalla stringa che indica il tipo (img, vid, aud) e il carattere ":";
		\item deve contenere il path del file;
		\item opzionalmente può contenere le coordinate X e Y del media, preceduti dai caratteri ":x." e ":y.".
	\end{itemize}
	Segue un esempio di QML con allegato per la domanda.
	\begin{verbatim}|q?open#t#Chi è l'autore di quest'opera?{img:laPass.png:x.5:y.5}#a#Monet#££#|\end{verbatim}
	
	\subsubsection{Struttura JSON}
	\label{qson}
	La struttura JSON delle domande deve seguire questa gerarchia:
	\begin{itemize}
		\item type;
		\item \begin{itemize}
			\item title/arrayTitle (se la domanda è singola o in piu parti), che contiene stringhe o JSON, il nome del campo può variare in base al tipo di domanda, quindi occorre progettare il parser e il generator tenendo conto del nome del campo utilizzato,
			\item correctAnswer/arrayAnswer (se la risposta è singola o in piu parti), che contiene stringhe o JSON, il nome del campo può variare in base al tipo di domanda, quindi occorre progettare il parser e il generator tenendo conto del nome del campo utilizzato,
		\end{itemize}
	\end{itemize}
	Campo questionAttachment/answerAttachment è opzionale e contiene eventuali allegati con la seguente struttura:
	\begin{itemize}
		\item type,
		\item path,
		\item x,
		\item y.
	\end{itemize}
	Segue un esempio per le stringhe utilizzate prima.
	\begin{verbatim}
	{type: 'trfs',
	 question: {
	            title: 'Parigi è in Germania',
	            correctAnswer: 'false'
	           }	
	}
	
	{type: 'mult',
	 question: {
	            title: 'Quali sono città europee?',
	            arrayAnswer: [ {textAnswer: 'Londra', isTrue: 'true'},
	                           {textAnswer: 'Mosca', isTrue: 'false'} ]
	           }	
	}
	\end{verbatim}
	
	\subsection{QMLAgent}
	Si trova al percorso \centra{app/controller/QMLAgent.js} 
	È il componente che si occupa di tradurre il QML in JSON, per poter poi creare effettivamente la domanda in modo visuale, e da JSON a QML per il salvataggio di una domanda nel database.
	Questo dopo aver rilevato il tipo di domanda chiama la funzione opportuna che genera il JSON o la stringa QML.
	Qualora si volesse implementare una nuova tipologia è quindi fondamentale seguire il seguente procedimento:
	\begin{itemize}
		\item scegliere una stringa univoca per identificare il tipo di domanda;
		\item inserire il case che riconosca tale stringa nello switch delle funzioni \texttt{generate} e \texttt{parse};
		\item aggiungere nel file la funzione che effettui il parsing della stringa QML a JSON, che rispetti le regole definite nella sottosezione \hyperref[qml]{Sintassi QML}.
		\item aggiungere nel file la funzione che effettui la generazione della stringa QML dal JSON, che rispetti le regole definite nella sottosezione \hyperref[qml]{Sintassi QML}.
		\item inserire la chiamata delle due funzioni nel case definito al punto 2.
	\end{itemize}
	scegliere una stringa univoca per identificare il tipo di domanda e inserire il caso che riconosca tale stringa.
	Successivamente è necessario creare una funzione di generazione e una di parsing per questo nuovo tipo, seguendo le regole definite in precedenza ed eventualmente creandone di nuove.
	Per il formato del JSON delle domande si veda la sottosezione \hyperref[qson]{Struttura JSON}.
	
	\subsection{CtrlQuestion}
	Si trova al percorso \centra{public/question\_creation/CtrlQuestion.js} 
	È il componente che si occupa di gestire la creazione di domande lato client.
	Qualora si volesse implementare una nuova tipologia è quindi fondamentale seguire il seguente procedimento:
	\begin{itemize}
		\item prima bisogna creare la factory della domanda che sarà il model della domanda lato client, e salvare il file js nella directory public/services;
		\item nelle dipendenze di CtrlQuestion aggiungere il nome del factory della domanda;
		\item assegnare allo \$scope una variabile che sarà del tipo della domanda che vogliamo creare;
		\item modificare lo switch del metodo \texttt{\$scope.createQuestion} aggiungendo il case relativo al nuovo tipo di domanda;
		\item nel corpo del case creare un metodo \texttt{save} che si occuperà solo del salvataggio della nuova domanda da salvare;
		\item definire il metodo \texttt{save} della domanda nello \$scope.
	\end{itemize}
	
	\subsection{CreateQuestionBase}
	Si trova al percorso \centra{public/question\_creation/CreateQuestionBase.html}
	È il componente che si occupa di mostrare la creazione dei vari tipi di domande.
	Qualora si volesse implementare una nuova tipologia è quindi fondamentale seguire il seguente procedimento:
	\begin{itemize}
		\item creare una pagina html con il nome della domanda che conterrà solo il codice per rappresentare a video la struttura della domanda, e salvare il file html nella directory public/question\_creation;
		\item creare una direttiva \gl{angular} di tipo elemento che esegue il fetch del path della pagina html precedentemente creata, e salvare il file js nella directory public/question\_creation;
		\item dentro createQuestionBase inserire tanti tag \texttt{src} quanti sono i file js appena creati (direttiva e factory);
		\item dentro la "variable part" di createQuestionBase aggiungere la direttiva relativa alla pagina html creata al punto 1.
	\end{itemize}
	
	\subsection{CtrlQuizExec}
	Si trova al percorso \centra{public/quiz\_execution/CtrlQuizExec.html}
	È il componente che si occupa di gestire lo svolgimento di quiz lato client.
	Qualora si volesse implementare una nuova tipologia è quindi fondamentale seguire il seguente procedimento:
	\begin{itemize}
		\item 
	\end{itemize}
	
	
	\subsection{Manuale Utente}
	Una volta inserito un nuovo tipo di domanda sarà opportuno inserire una sua descrizione adeguata in appendice al \doc{Manuale Utente}.
	
	\newpage
	\section{Account SuperAdmin}
	Una volta installato il prodotto viene creato di default un account SuperAdmin che permette di creare enti e l'account responsabile associato. 
	\\ Le credenziali di login per SuperAdmin sono:
	\begin{description}
		\item[Username] SuperAdmin,
		\item[Password] la password sarà generata random alla prima installazione e sarà reperibile nel file \texttt{config/admin.txt}.
	\end{description}
	Una volta loggato con le credenziali si troverà un form speciale che permette l'inserimento nel database di enti, account responsabili associati e una pagina che permette di vedere le varie segnalazioni e richieste inoltrate dagli utenti che utilizzano il form "Segnalazione e richieste". Seguono degli screenshot di queste pagine.
	
	\begin{figure}[!h]
		\centering
		\includegraphics[scale=0.5]{Img/logo.png}
		\caption{Schermata di creazione enti e nomina responsabili}
	\end{figure}
	
	\begin{figure}[!h]
		\centering
		\includegraphics[scale=0.5]{Img/logo.png}
		\caption{Schermata di segnalazioni e richieste}
	\end{figure}
	
	ATTENZIONE: utilizzare l'account SuperAdmin è l'unico metodo per creare enti e nominare responsabili, e quindi almeno un accesso è obbligatorio per rendere fruibile al pubblico il prodotto dopo il primo avvio (dopo che qualche utente avrà inoltrato richiesta di creazione ente).
	
	\section{Segnalazione di problemi d'uso}
	Nel caso si verificassero dei problemi d'uso di entità più grave è possibile contattare il team Vault-Tech tramite email \centra{vaulttech.swe@gmail.com}
	
	
	\newpage
	\appendix
	
	\section{Glossario}
	\label{gl} 
	
	\subsection{Angular.js}
	AngularJS (o semplicemente Angular o Angular.js) è un framework web open source principalmente sviluppato da Google e dalla comunità di sviluppatori individuali che ruotano intorno al framework nato per affrontare le molte difficoltà incontrate nello sviluppo di applicazioni singola pagina. Ha l'obiettivo di semplificare lo sviluppo e il test di questa tipologia di applicazioni fornendo un framework lato client con architettura MVC (Model View Controller) e Model–View–ViewModel (MVVM) insieme a componenti comunemente usate nelle applicazioni RIA.
	
	\subsection{Browser}
	Programma che consente di usufruire dei servizi di connettività in Internet, o di una rete di computer.
	
	\subsection{CSS3}
	Il CSS (Cascading Style Sheets) è un linguaggio usato per definire la formattazione di documenti HTML, XHTML e XML. CSS3 è l’ultimo standard approvato al momento.
	
	\subsection{Git}
	Sistema software di controllo di versione distribuito, utilizzabile da riga di comando.
	
	\subsection{HTML5}
	L’HTML5 è un linguaggio di markup per la strutturazione delle pagine web divenuto standard W3C nell’ottobre 2014.
	
	\subsection{JavaScript}
	Linguaggio di programmazione interpretato, generalmente utilizzato nella gestione degli eventi.
	Implementa un paradigma basato sia sugli oggetti che sulla programmazione funzionale.
	
	\subsection{LTS}
	 Supporto a lungo termine (Long Term Support), usato per programmi che indicano che la versione di quel programma fornirà un supporto per un periodo più lungo di tempo. Uno ogni quattro rilasci di Ubuntu è LTS.
	
	\subsection{MongoDB}
	MongoDB è un Data Base Management System non relazionale, orientato ai documenti. Classificato come un database di tipo NoSQL.
	
	\subsection{Node.js}
	Node.js è un framework event-driven per il motore JavaScript V8, su piattaforme UNIX like, si tratta cioè di un framework relativo all'utilizzo server-side di JavaScript.
	
	\subsection{Npm}
	Npm è il package manager di default per Node.js. Permette di installare moduli extra per Node.js da riga di comando.
	
	\subsection{Repository}
	È un ambiente di un sistema informativo, in cui vengono gestiti i dati, attraverso tabelle relazionali. In questo caso il sistema informativo è gestito con Git.
	
	\subsection{Server}
	Un server in informatica è un componente o sottosistema informatico di elaborazione e gestione del traffico di informazioni che fornisce, a livello logico e fisico, un qualunque tipo di servizio ad altre componenti (tipicamente chiamate clients, cioè clienti) che ne fanno richiesta attraverso una rete di computer, all'interno di un sistema informatico o anche direttamente in locale su un computer.
	
	\subsection{Software}
	Un software, in informatica, è l’informazione o le informazioni utilizzate da uno o più sistemi informatici
	e memorizzate su uno o più supporti informatici. Tali informazioni possono essere quindi rappresentate da uno o più programmi, oppure da uno o più dati, oppure da una combinazione
	delle due.
	
	\subsection{Web}
	Il Web, abbreviazione di World Wide Web, è uno dei principali servizi di Internet che permette di
	navigare e usufruire di un insieme vastissimo di contenuti (multimediali e non) collegati tra loro
	attraverso legami (link), e di ulteriori servizi accessibili a tutti o ad una parte selezionata degli
	utenti di Internet.
	
\end{document}
